\documentclass[conference,10pt]{IEEEtran}

\usepackage[T1]{fontenc}
\usepackage{microtype}
\usepackage{amsmath}
\usepackage{amssymb}
\usepackage{booktabs}
\usepackage{graphicx}
\usepackage{multirow}
\usepackage{xcolor}
\usepackage[hidelinks]{hyperref}
\usepackage{url}

\newcommand{\system}{\textsc{FamilyPhase}}
\newcommand{\draftwarning}{%
  \begin{center}
  \fbox{\parbox{0.94\columnwidth}{\footnotesize\textbf{Internal evidence-audit draft---do not submit.}
  The final multicore files are not mounted in the current environment. Numbers marked
  ``pilot'' come from the older five-workload validation set and are included to test
  the analysis and paper structure, not to replace the required final rerun.}}
  \end{center}}

\title{One Counter per Resource Family: Compact Online Multicore Phase Prediction}
\author{\IEEEauthorblockN{Anonymous Author(s)}}

\begin{document}
\maketitle

\begin{abstract}
Hardware performance counters can expose workload phases, but a full counter vector
often exceeds the physical PMU slots and a high-accuracy phase predictor can obtain a
misleading result by merely repeating an oracle phase label. We study a stricter
interface: an offline, train-only full-counter teacher defines a global pressure phase,
while an online predictor observes one validation-selected counter state from each
resource family. We separate oracle diagnostics from deployable baselines, split and
resample complete executions, and include the discretizer and history state in the
storage budget. In an audited pilot of 60 PARSEC executions, a depth-six tree over five
20-interval histories obtains 98.21\% held-out accuracy (95\% execution-group bootstrap
CI: 97.64--98.71\%) using an estimated 201.25 bytes. Against an equally deployable tree
using only the current counter states, history improves accuracy by 0.60 percentage
points (paired 95\% CI: 0.29--0.92) and transition accuracy by 12.95 points, although
the latter CI includes zero. The audit also finds material multiplexing variation in the pilot
collection, showing why enabled-time and workload-disjoint results must accompany
accuracy. These results motivate a measurement-quality-first evaluation of compact
phase interfaces rather than an uncalibrated claim of end-to-end hardware savings.
\end{abstract}

\draftwarning

\section{Introduction}

An online hardware controller needs a compact description of what the workload will
do next. Raw performance-monitoring-unit (PMU) events are natural signals, yet three
problems make a phase interface difficult to evaluate. First, processors expose more
useful events than can be scheduled simultaneously; multiplexing can produce noisy or
missing 10~ms samples. Second, adjacent phase labels are highly persistent, so a model
can report high accuracy by repeating the current phase while detecting no changes.
Third, randomly splitting intervals leaks nearly identical samples from one execution
into both training and test sets. These problems matter before the choice of model.

Prior work established code- and counter-based program phases for representative
simulation, runtime prediction, cache adaptation, and multicore management
\cite{sherwood2001bbda,sherwood2002simpoint,sherwood2003phase,perelman2006parallel,
nomani2015predicting,lozano2022phase}. Rich offline signatures are valuable for phase
discovery but are not automatically viable as a short-latency hardware interface.
Conversely, a tiny model is not useful if it relies on an oracle current label or a
counter set that the PMU rarely schedules.

We investigate whether a global phase discovered from a broad safe-counter vector can
be approximated online by one counter per resource family. The families correspond to
L1, L2, LLC, branch/control, core/FP, and memory/off-core behavior. A train-only
$k$-means teacher creates ordered pressure labels. Validation ablation freezes one
event per available family. One-dimensional train-only clustering converts each event
to a two-bit state, and a shallow tree predicts the next global label from the last
$H$ states. We call this pipeline \system.

The paper makes three methodological contributions.

\begin{itemize}
  \item We define a reduced resource-family PMU interface while keeping the target
  explicit: every student predicts one \emph{global full-counter teacher label}; the
  per-family states are inputs, not separate phase ground truths.
  \item We introduce an evaluation protocol that groups co-running processes and
  repetitions before splitting, labels persistence/Markov models that consume the
  teacher's current phase as oracle diagnostics, and compares history against a
  deployable current-state-only tree.
  \item We make collection quality and implementation cost first-class outputs:
  execution-group confidence intervals, $k$/seed stability, PMU enabled time, tree
  state, discretization thresholds, and history storage are reported together.
\end{itemize}

The audited pilot suggests that compact histories contain useful transition
information, but it also exposes evidence that must be corrected before a full claim:
only nine runs are held out in the primary split, $k=2$ is more stable than the chosen
$k=3$, and several events have poor enabled time. We report these limitations rather
than converting them into a proxy energy result.

\section{Background and Related Work}

\subsection{Program phases}

BBV and SimPoint methods cluster intervals with similar basic-block distributions and
use representative intervals to reduce simulation cost
\cite{sherwood2001bbda,sherwood2002simpoint,hamerly2005simpoint3}. Online phase trackers
replace long profiling signatures with compact histories and transition models
\cite{sherwood2003phase}. The definition is task dependent: working sets are useful for
cache allocation, while PMU events expose resource pressure with lower software
instrumentation cost \cite{dhodapkar2002managing,dhodapkar2003comparing}.

Multicore execution complicates phase ownership. A global label represents a package
or workload group cheaply, a local label represents each core independently, and a
distributed predictor combines local and cross-core context
\cite{perelman2006parallel,chang2013sampling,lozano2022phase}. \system uses a global
target because package-level cache, bandwidth, thermal, and scheduling policies need a
common state. Resource-family inputs describe which subsystems provide evidence for
that target; we do not call them family-specific target phases.

\subsection{Counter selection and phase-aware control}

Counter-based predictors have been applied to power/performance modeling, scheduling,
and cache-demand estimation \cite{kim2017p4,nomani2015predicting,ziedan2016runtime}.
Phase-aware policies have also controlled caches, dormant structures, DVFS, and power
balancing \cite{balasubramonian2000memory,albonesi1999selective,
buyuktosunoglu2000drowsy,isci2003runtime,merkel2006balancing}. These works motivate the
resource-family organization. Our focus is narrower: how much phase fidelity survives
an aggressively reduced observable interface, and how should that trade-off be
measured without leaking executions or granting a baseline oracle labels?

PMU event meaning and availability vary across microarchitectures. Multiplexed events
are scaled by Linux \texttt{perf}, but enabled time still controls measurement quality
and should be reported \cite{intel2024optimization,intel2024sdm,linux2024perf,
perfwiki2024}. This motivates separating teacher collection, reduced online collection,
and model evaluation instead of counting only final model parameters.

\section{Resource-Family Phase Interface}
\label{sec:design}

\subsection{Train-only global teacher}

Let $x_{r,t}\in\mathbb{R}^{d}$ be the safe PMU vector at interval $t$ of run $r$.
Timing-derived features such as cycles, IPC/CPI, elapsed time, and requested frequency
are excluded so that a future policy does not trivially feed its action back into the
label. Each nonnegative count is transformed by $\log(1+x)$, missing values are filled
with training medians, and features are standardized with training statistics.

The teacher fits $k$-means on training rows only:
\begin{equation}
  c_t=\arg\min_{j\in\{0,\ldots,k-1\}}\|\tilde{x}_{r,t}-\mu_j\|_2^2.
\end{equation}
Validation and test rows use the frozen scaler and nearest centroid. Anonymous cluster
IDs are ordered by an increasing pressure score over available LLC-miss, memory-
bandwidth, branch-misprediction, and stall indicators. For $k=3$, the labels are
reported as low, moderate, and high pressure. The labels are descriptive operating
regions, not an external ground truth for power or bottlenecks.

\subsection{Selecting and discretizing the online interface}

Candidate events are partitioned into resource families. Near-duplicate events are
identified on training data using Spearman correlation. Singleton decision-tree
ablations are scored on validation accuracy and high-pressure recall, and one
representative $e_f^\star$ is frozen per available family $f$. The test split never
selects an event.

For the online representation, a one-dimensional three-cluster model is fit to
$e_f^\star$ on training rows. Its ordered state
$s_{f,t}\in\{0,1,2\}$ needs two bits. Two 16-bit boundaries per family approximate the
nearest-centroid decision. The cross-family history vector is
\begin{equation}
  u_t=[s_{1,t-H+1},\ldots,s_{F,t-H+1},\ldots,s_{F,t}],
\end{equation}
and the deployable predictor learns $h(u_t)\rightarrow c_{t+1}$. A depth-limited CART
tree is used because inference consists of comparisons and table-like child traversal
\cite{breiman1984cart,quinlan1993c45}.

\subsection{Baselines and their information sets}

Table~\ref{tab:baselines} prevents an easy but consequential category error. Persistence,
state-conditioned majority, Markov, run-length Markov, and the duration approximation
consume the teacher's true $c_t$. Because the full teacher is offline, these methods
are oracle diagnostics unless a separate online phase estimator is added. They measure
phase persistence but are not fair hardware implementations. The majority predictor
uses only the training label prior. The fair history ablation is a tree trained on the
current vector $[s_{1,t},\ldots,s_{F,t}]$ with the same depth/leaf limits.

\begin{table}[t]
\centering
\caption{Model information available at prediction time.}
\label{tab:baselines}
\footnotesize
\begin{tabular}{p{0.31\columnwidth}p{0.38\columnwidth}p{0.13\columnwidth}}
\toprule
Model & Online input & Deploy. \\
\midrule
Majority & training prior & yes \\
Teacher persistence/Markov & true teacher $c_t$ & no (oracle) \\
Current-state tree & selected $s_{1..F,t}$ & yes \\
Single-family history tree & selected $s_{f,t-H+1:t}$ & yes \\
Cross-family history tree & selected $s_{1..F,t-H+1:t}$ & yes \\
\bottomrule
\end{tabular}
\end{table}

\subsection{Analytical implementation budget}

For a tree with $I$ internal nodes, $L$ leaves, $D$ input features, and $N=I+L$
nodes, the model estimate is
\begin{equation}
 B_{tree}=\frac{I(\lceil\log_2D\rceil+16+2\lceil\log_2N\rceil)+2L}{8}.
\end{equation}
The full online estimate adds $\lceil2FH/8\rceil$ history bytes and $4F$ bytes for two
16-bit discretizer boundaries per family. This excludes PMU register hardware already
present in the processor and is an analytical storage estimate, not synthesis area.

\section{Experimental Methodology}

\subsection{Audited pilot}

The currently reproducible pilot uses five PARSEC workloads---blackscholes, bodytrack,
canneal, fluidanimate, and freqmine---with 1, 2, 4, and 8 threads and three repetitions
per configuration (60 runs) \cite{bienia2008parsec}. Task-local counters were sampled
at 10~ms on a dual-socket Intel Xeon E5-2680~v2 system with 20 physical cores and SMT
disabled, using Linux \texttt{perf} 5.14. The safe merged table contains 9,409 interval
rows and eight usable counter columns. L2 events were unavailable, leaving five online
families. This dataset is a method-validation pilot, not the intended final multicore
co-runner study.

The primary split groups all repetitions of each workload/thread configuration before
assigning 70\%/15\%/15\% to train/validation/test. There are 5,642 valid history windows
for training and 1,341 for test; only nine complete runs occur in the test split.
Leave-one-workload-out results are generated separately. All interval metrics are
accompanied by a cluster bootstrap that samples complete run or concurrent-group IDs.

\subsection{Final multicore protocol}

The final study uses three regimes: one multithreaded process, two single-threaded
co-runners, and two four-thread co-runners. For seven workloads the collector enumerates
all $\binom{7}{2}=21$ unordered cross-workload pairs. Both processes in a concurrent
group share a split, and any workload holdout moves every group containing that
workload to test. This corrects an earlier exploratory design that used only four fixed
pairs and could split partners independently. No result from that old split is used in
this draft.

\subsection{Metrics and uncertainty}

We report accuracy, macro F1, balanced accuracy, high-pressure recall, stable-case
accuracy, and transition-case accuracy. A transition case satisfies $c_{t+1}\ne c_t$.
Final results will additionally report transition precision/recall/F1, false alarms,
and transition detection delay. Confidence intervals use 10,000 execution-group
bootstrap resamples. Model comparisons use paired resampling of the same groups; a
difference is not described as reliable when its 95\% interval crosses zero.

Clustering robustness is evaluated for $k\in\{2,3,4,5\}$ and five initializations using
pairwise adjusted Rand index (ARI), train distance, centroid separation, and minimum
cluster share. History, tree depth, and sampling-interval sensitivity are also required
for the final study.

\section{Audited Pilot Results}

\subsection{Does history help a deployable predictor?}

Table~\ref{tab:pilot-main} reports one copy of each cross-family result. The exporter
uses the family label \texttt{\_\_all\_families\_\_} for these rows because the global
target and cross-family features must not be repeated and averaged as independent
observations. The oracle persistence diagnostic
reaches 97.69\% accuracy because only 31 of 1,341 test windows transition, but it has
zero transition accuracy. Its high accuracy is evidence of label persistence, not a
deployable win.

The current-state tree recovers 97.61\% accuracy and correctly predicts 12.9\% of
transitions. The 20-interval history tree improves accuracy to 98.21\% and transition
accuracy to 25.8\%. The execution-group CIs are wider than an interval bootstrap would
suggest. A paired comparison finds a 0.60-point accuracy improvement (95\% CI:
0.29--0.92; bootstrap probability of a positive difference 99.99\%). The transition
improvement is 12.95 points, but its 95\% CI is 0.0--26.92 points. Thus this pilot
supports the accuracy effect but is underpowered for a strong transition claim.
As a binary change detector, the history tree has 90.0\% precision, 29.0\% recall,
and 43.9\% F1 (9 true changes and one false alarm); the current-state tree has 54.5\%
precision, 19.4\% recall, and 28.6\% F1 (6 true changes and five false alarms).

\begin{table}[t]
\centering
\caption{Pilot configuration-group holdout. CIs resample nine test runs. Oracle
persistence consumes the offline teacher phase and is not deployable.}
\label{tab:pilot-main}
\footnotesize
\begin{tabular}{lccc}
\toprule
Model & Accuracy (95\% CI) & Macro F1 & Trans. acc. \\
\midrule
Majority & 77.03 (70.77--82.69) & 29.01 & 48.4 \\
Oracle persistence & 97.69 (97.20--98.14) & 95.97 & 0.0 \\
Current-state tree & 97.61 (96.84--98.30) & 95.45 & 12.9 \\
20-history tree & \textbf{98.21} (97.64--98.71) & \textbf{96.88} & \textbf{25.8} \\
\bottomrule
\end{tabular}
\end{table}

The majority transition accuracy is not a contradiction: with few transitions, a
constant phase can happen to match one transition destination while failing most
stable regions. This is why neither overall accuracy nor transition accuracy alone is
sufficient. Macro F1 and the confusion matrix must accompany both.

\subsection{Does the interface generalize to an unseen workload?}

Table~\ref{tab:pilot-loo} holds out every run of one workload at a time. Across the
five folds, the history tree averages 94.60\% accuracy and 43.9\% transition accuracy,
compared with 93.01\% and 22.5\% for the current-state tree. Oracle persistence averages
94.03\% accuracy and zero transition accuracy. The result is not uniformly easy:
bodytrack reduces history-tree accuracy to 82.24\%. This failure case should be
explained with its phase shares and confusion matrix rather than hidden by the mean.
Paired execution-group bootstraps place the transition-accuracy improvement above zero
in every fold (the smallest lower bound is 5.9 points). Accuracy improves reliably in
four folds; canneal's 0.07-point difference has a 95\% CI of $-0.38$--0.63 points.
Because folds reuse training workloads, their arithmetic mean is descriptive rather
than five independent replications.

\begin{table}[t]
\centering
\caption{Pilot leave-one-workload-out results (percent).}
\label{tab:pilot-loo}
\scriptsize
\setlength{\tabcolsep}{3pt}
\begin{tabular}{lrrrr}
\toprule
Holdout & Cur. acc. & Hist. acc. & Cur. trans. & Hist. trans. \\
\midrule
blackscholes & 97.68 & 98.24 & 21.43 & 40.48 \\
bodytrack & 76.47 & 82.24 & 21.04 & 47.26 \\
canneal & 98.22 & 98.28 & 25.71 & 45.71 \\
fluidanimate & 95.77 & 96.69 & 20.27 & 50.00 \\
freqmine & 96.94 & 97.55 & 24.14 & 36.21 \\
\midrule
Mean & 93.01 & \textbf{94.60} & 22.52 & \textbf{43.93} \\
\bottomrule
\end{tabular}
\end{table}

\subsection{What does the compact interface cost?}

The current-state tree uses an estimated 64.88~B including five discretizers and a
two-byte current-state register. The cross-family history tree uses 201.25~B including
25~B of two-bit history and 20~B of thresholds. The increase is small in absolute terms,
but it should not be advertised as RTL area without synthesis. Software inference
latency in the Python implementation is approximately 52~$\mu$s per prediction and is
not a cycle-level hardware estimate.

\subsection{Are three teacher states stable?}

Table~\ref{tab:k-stability} shows that the pilot does not uniquely justify $k=3$.
Three states have mean pairwise ARI 0.927, but the smallest cluster averages only 2.1\%
of rows. Two states are almost invariant to initialization and more balanced. Four and
five states improve within-cluster distance but reduce seed stability and create even
smaller clusters. A final paper may retain $k=3$ for low/moderate/high semantics only
if the larger multicore study confirms stability and utility; otherwise $k$ should be
selected on train/validation evidence.

\begin{table}[t]
\centering
\caption{Pilot train-only clustering sensitivity over five seeds.}
\label{tab:k-stability}
\footnotesize
\begin{tabular}{ccccc}
\toprule
$k$ & Mean ARI & Min ARI & Train dist. & Min share \\
\midrule
2 & 1.000 & 0.999 & 1.418 & 15.8\% \\
3 & 0.927 & 0.817 & 1.276 & 2.1\% \\
4 & 0.845 & 0.742 & 1.136 & 1.8\% \\
5 & 0.734 & 0.511 & 0.990 & 1.8\% \\
\bottomrule
\end{tabular}
\end{table}

\subsection{PMU scheduling is a first-order validity issue}

Reprocessing exactly the 60 raw pilot runs preserves \texttt{perf}'s enabled percentage.
For the eight safe events used by the teacher, mean enabled time is 90.6--91.0\%, but
the fifth percentile is only 64--66\% and 32.4--33.7\% of samples fall below 90\%.
Rare samples reach zero enabled time. This variation can distort both the teacher and
the selected interface at 10~ms granularity. Accordingly, the final experiment must
compare the full profile against the frozen reduced profile and report enabled-time
improvement and runtime overhead. The pilot accuracy cannot establish that benefit by
itself.

\section{Discussion}

\subsection{What the result supports}

The pilot supports a limited conclusion: within seen workload families and a grouped
configuration split, five compact current states approximate the next global teacher
phase, and a short history provides a small paired accuracy improvement. The analytical
state is hundreds rather than tens of thousands of bytes. These properties make the
interface plausible for a firmware or hardware prototype.

It does not establish energy savings, cache savings, or optimal control. The prior
draft mapped labels to arbitrary L1 power factors; that proxy is removed because no
cache size, miss latency, energy model, or slowdown was calibrated. A control claim
requires Linux \texttt{resctrl}, a cache simulator, DVFS measurements, or another real
actuator with performance and energy accounting.

\subsection{Threats to validity}

\textbf{Construct validity.} The teacher labels are clusters in the chosen counter
space, not independently measured bottlenecks. Ordering them by pressure aids
interpretation but does not prove that every high label needs the same action.

\textbf{Internal validity.} PMU multiplexing, frequency variation, OS noise, NUMA
placement, affinity, warm-up, and interval alignment can change counts. The final
artifact must report event enabled time, governor/turbo state, kernel and \texttt{perf}
versions, exact event encodings, pinning, socket placement, unsuccessful samples, and
run randomization.

\textbf{Statistical validity.} Intervals from one run are correlated. We therefore
split and bootstrap runs/concurrent groups, but nine pilot test runs remain too few for
a transition conclusion. Multiple seeds and paired group resampling reduce, but do not
eliminate, uncertainty.

\textbf{External validity.} The pilot covers five PARSEC programs on one Ivy Bridge
server with task-local events. Results may not transfer to different PMUs, production
services, SMT, other sockets, or unseen co-runners. Workload-disjoint, pair-disjoint,
and preferably second-generation CPU experiments are required.

\textbf{Implementation validity.} Byte and comparison counts are analytical. They
exclude counter read energy, control integration, timing closure, and actuator cost.
RTL synthesis or carefully qualified firmware measurements are needed before calling
the predictor implemented hardware.

\section{Conclusion}

Compact phase prediction should be evaluated as a measurement interface, not only as
a classifier. \system separates an offline global teacher from one-counter-per-family
online signals, distinguishes oracle persistence from deployable baselines, and reports
execution-group uncertainty and complete interface storage. An audited pilot finds a
small but reliable accuracy gain from history and a promising, not yet conclusive,
transition gain. It also reveals material PMU enabled-time variation and $k$ sensitivity. Those
findings define the evidence required for the final multicore study: all co-runner
pairs, reduced-profile overhead, workload-disjoint tests, robust phase selection, and
no proxy energy claims.

\bibliographystyle{IEEEtran}
\bibliography{references}

\end{document}
